%% kapa: exam=1 round=25 year=2014 period=1 form=A subject=economics subject_ko=경제원론 seq=4 number=44
%% source: Q-Net 감정평가사 기출문제 (cst00309) artlSeq=5101207 file=extracted/1008/5101207/2014년 제25회 감정평가사 1차시험 기출문제/1교시(A형).hwp
%% conversion: kordoc 4.24.51 -> tools/gen1.py (auto); review=proofread-sonnet; flags=none
\begin{question}{44}
효용을 극대화하는 소비자 甲이 $X$재와 $Y$재 두 재화만 소비한다. $X$재의 가격이 하락할 때 다음 설명 중 옳은 것을 모두 고른 것은?
\begin{framedtext}
ㄱ. $Y$재가 열등재이면, $Y$재의 수요는 감소한다. \par
ㄴ. $X$재와 $Y$재가 모두 필수재이면, $X$재에 대한 지출액은 감소한다. \par
ㄷ. $X$재와 $Y$재가 모두 열등재이면, $Y$재의 수요는 증가한다. \par
\end{framedtext}
\begin{choices}
  \item ㄴ
  \item ㄱ, ㄴ
  \item ㄱ, ㄷ
  \item ㄴ, ㄷ
  \item ㄱ, ㄴ, ㄷ
\end{choices}
\end{question}
\answerall
\answernote{정답표에 1,2,3,4,5 모두 정답으로 표기(전원 정답)}
